% This file is part of the modular source for Stratified Relational Workspace Theory.
% !TeX root = ../Stratified_Relational_Workspace_Theory.tex
\theorypart{IV}{Approximate Workspaces}

The exact theory asks whether local core states agree perfectly under every
transfer. Approximate theory asks three more realistic questions:
\begin{enumerate}
\item how strongly a candidate state violates edgewise compatibility;
\item whether a low-defect subspace is spectrally separated from incoherent
 directions; and
\item how that subspace changes when transfers, metrics, or operations are
 perturbed.
\end{enumerate}
Throughout Parts~IV--V, set
\[
\mathcal C_V=\bigoplus_{i\in V}C_i,
\qquad
\mathcal C_E=\bigoplus_{e\in E}C_{t(e)},
\]
with their Hilbert direct-sum inner products, and retain strictly positive
weights \(w_e\).

\topic{11}{Transfer Discrepancy Operator}

\begin{definition}[Transfer discrepancy]\label{def:transfer-discrepancy}
The transfer discrepancy operator is
\[
D_H:\mathcal C_V\longrightarrow\mathcal C_E,
\qquad
(D_Hc)_e
=
\sqrt{w_e}\bigl(c_{t(e)}-H_ec_{s(e)}\bigr).
\]
Its squared norm is the coherence energy
\[
\boxed{
\mathcal E_H(c)
=
\|D_Hc\|^2
=
\sum_{e\in E}w_e
\left\|c_{t(e)}-H_ec_{s(e)}\right\|^2.
}
\]
The transfer Laplacian is
\[
\boxed{L_H=D_H^*D_H.}
\]
\end{definition}

The operator \(D_H\) is an incidence-type derivative twisted by the edge
transfers. It compares a state already present at the target with the state
predicted there by translation from the source. A direction can therefore
have low energy even when its coordinates differ substantially across nodes,
provided those coordinates differ exactly as the transfer maps prescribe.

For \(z=(z_e)_{e\in E}\in\mathcal C_E\), the adjoint is
\[
(D_H^*z)_i
=
\sum_{e:t(e)=i}\sqrt{w_e}\,z_e
-
\sum_{e:s(e)=i}\sqrt{w_e}\,H_e^*z_e.
\]
Consequently,
\[
\boxed{
\begin{aligned}
(L_Hc)_i
={}&
\sum_{e:t(e)=i}w_e\bigl(c_i-H_ec_{s(e)}\bigr)\\
&+
\sum_{e:s(e)=i}w_eH_e^*
\bigl(H_ec_i-c_{t(e)}\bigr).
\end{aligned}
}
\]
This formula makes the two contributions at a vertex explicit: incoming
prediction error and adjoint-propagated outgoing error.

\begin{proposition}[Positivity and variational identity]\label{prop:laplacian-positive}
The transfer Laplacian is self-adjoint and positive semidefinite. For every
\(c\in\mathcal C_V\),
\[
\langle c,L_Hc\rangle=\mathcal E_H(c)\geq0.
\]
\end{proposition}

\begin{proof}
Self-adjointness and positivity follow from \(L_H=D_H^*D_H\). The displayed
identity is
\[
\langle c,D_H^*D_Hc\rangle
=
\langle D_Hc,D_Hc\rangle.
\]
\end{proof}

\begin{remark}[Weight dependence]
Replacing \(w_e\) by another positive family does not alter \(\ker D_H\), but
it generally changes every positive eigenvalue and therefore changes the
approximate geometry. Edge weights are part of the metric model, not part of
the exact compatibility relation.
\end{remark}

\topic{12}{Exact Kernel Theorem}

\begin{theorem}[Transfer-Laplacian kernel theorem]\label{thm:laplacian-kernel}
For an arbitrary finite directed core diagram, without assuming invertibility,
\[
\boxed{
\ker L_H=\ker D_H=\Gamma(H).
}
\]
Thus the exact compatible sections are precisely the zero-energy modes of the
transfer Laplacian.
\end{theorem}

\begin{proof}
If \(D_Hc=0\), then \(L_Hc=D_H^*D_Hc=0\). Conversely, if \(L_Hc=0\), then
\[
0=\langle c,L_Hc\rangle=\|D_Hc\|^2,
\]
so \(D_Hc=0\). Finally, \(D_Hc=0\) if and only if
\[
c_{t(e)}=H_ec_{s(e)}
\]
for every edge, which is exactly the definition of \(\Gamma(H)\).
\end{proof}

\begin{corollary}[Finite rank computation]\label{cor:workspace-rank-D}
Let
\[
n_C=\sum_{i\in V}\dim C_i.
\]
Then
\[
\boxed{
\dim\Gamma(H)=n_C-\operatorname{rank}(D_H).
}
\]
Hence the exact workspace dimension can be obtained by a single rank or
singular-value computation, even when holonomy is unavailable because some
transfers are noninvertible.
\end{corollary}

\subtopic{12.1 Relation to the holonomy computation}

In the local-system regime, Theorems~\ref{thm:limit-holonomy} and
\ref{thm:laplacian-kernel} give two equivalent constructions:
\[
\Gamma(H)
\cong
C_o^{\Hol_o}
\qquad\text{and}\qquad
\Gamma(H)=\ker L_H.
\]
The first reduces the problem to cycle invariants in a root fiber. The second
works directly on all vertex coordinates and extends to noninvertible diagrams.
The two computations therefore serve different practical regimes rather than
competing definitions.

\topic{13}{Approximate Structural Workspace}

Let
\[
0\leq\lambda_1\leq\lambda_2\leq\cdots\leq\lambda_N,
\qquad
N=\dim\mathcal C_V,
\]
be the eigenvalues of \(L_H\), repeated with multiplicity, and let
\(P_I(L_H)\) denote the spectral projector associated with an interval \(I\).

\begin{definition}[Threshold workspace]\label{def:threshold-workspace}
For \(\varepsilon\geq0\), define
\[
\boxed{
\Gamma_\varepsilon(H)
=
\ran P_{[0,\varepsilon]}(L_H).
}
\]
Its dimension is
\[
d_\varepsilon
=
\#\{k:\lambda_k\leq\varepsilon\}.
\]
When \(0<d_\varepsilon<N\), define the boundary gap
\[
\boxed{
g_\varepsilon
=
\lambda_{d_\varepsilon+1}-\lambda_{d_\varepsilon}.
}
\]
\end{definition}

The threshold \(\varepsilon\) declares which discrepancy scale is treated as
structurally negligible. The dimension \(d_\varepsilon\) alone is not a
robust certificate: a cluster is stable only when it is separated by a
positive boundary gap.

An equivalent index-based description fixes \(d\) and writes
\[
W_d=\ran P_d,
\qquad
P_d=P_{\{\lambda_1,\ldots,\lambda_d\}}(L_H),
\]
with cluster gap
\[
g_d=\lambda_{d+1}-\lambda_d>0.
\]
The index-based form is convenient for perturbation theorems because the
workspace rank is declared explicitly.

\begin{proposition}[Min--max characterization]\label{prop:minmax}
For \(1\leq k\leq N\),
\[
\lambda_k
=
\min_{\substack{S\subseteq\mathcal C_V\\\dim S=k}}
\ \max_{\substack{c\in S\\\|c\|=1}}
\mathcal E_H(c).
\]
Consequently, the span of the first \(k\) eigenvectors minimizes the worst-case
coherence energy among all \(k\)-dimensional candidate subspaces.
\end{proposition}

\begin{proof}
Expand every vector in an orthonormal eigenbasis of \(L_H\). The span of the
first \(k\) eigenvectors has maximal Rayleigh quotient \(\lambda_k\). Every
\(k\)-dimensional subspace intersects the orthogonal complement of the first
\(k-1\) eigenvectors nontrivially, and any unit vector in that intersection has
Rayleigh quotient at least \(\lambda_k\).
\end{proof}

\subtopic{13.1 Exact, soft, and spurious low modes}

The spectrum supports three distinct interpretations:
\[
\begin{array}{rcl}
\lambda_k=0
&\Longleftrightarrow&
\text{exact compatible direction},\\
0<\lambda_k\ll g_d
&\Longrightarrow&
\text{approximately transported direction},\\
\lambda_k\approx\lambda_{k+1}
&\Longrightarrow&
\text{rank selection is unstable at that boundary}.
\end{array}
\]
A low eigenvalue can still be scientifically spurious if it arises from a
poorly chosen weight, a nearly disconnected graph, or a noncausal encoding.
The causal and null gates in later parts remain necessary.

\topic{14}{Normalized Coherence Defect}

For \(c\neq0\), define the Rayleigh defect
\[
\delta_{\mathrm{coh}}(c)
=
\frac{\mathcal E_H(c)}{\|c\|^2}.
\]
For a nonzero subspace \(W\subseteq\mathcal C_V\), define
\[
\delta_{\mathrm{coh}}(W)
=
\sup_{\substack{c\in W\\\|c\|=1}}
\mathcal E_H(c).
\]
If \(W\) is \(L_H\)-invariant, then
\[
\delta_{\mathrm{coh}}(W)
=
\|L_H|_W\|.
\]
In particular,
\[
\delta_{\mathrm{coh}}(\Gamma_\varepsilon(H))
\leq\varepsilon.
\]

\begin{theorem}[Coherence-to-core distance estimate]\label{thm:poincare-workspace}
Let
\[
W_0=\Gamma(H)=\ker L_H,
\qquad
r=\dim W_0,
\]
and suppose \(r<N\). Set
\[
\gamma_H=\lambda_{r+1}>0.
\]
Then every \(c\in\mathcal C_V\) satisfies
\[
\boxed{
\operatorname{dist}(c,W_0)^2
\leq
\frac{\mathcal E_H(c)}{\gamma_H}.
}
\]
Moreover,
\[
\mathcal E_H(c)
\leq
\|L_H\|\operatorname{dist}(c,W_0)^2.
\]
\end{theorem}

\begin{proof}
Write \(c=P_0c+Q_0c\), where \(P_0\) projects onto \(W_0\) and
\(Q_0=I-P_0\). Since \(L_HP_0=0\),
\[
\mathcal E_H(c)
=
\langle Q_0c,L_HQ_0c\rangle.
\]
The spectrum of \(L_H\) on \(W_0^\perp\) lies in
\([\gamma_H,\|L_H\|]\). Hence
\[
\gamma_H\|Q_0c\|^2
\leq
\mathcal E_H(c)
\leq
\|L_H\|\|Q_0c\|^2.
\]
Finally, \(\|Q_0c\|=\operatorname{dist}(c,W_0)\).
\end{proof}

This estimate is the exact sense in which a spectral gap turns small edgewise
disagreement into closeness to a genuinely compatible global section.

\subtopic{14.1 Two-node consistency cycle}

Let \(C_1=C_2=\mathbb R\), use unit weights, and take two directed transfers
\[
H_{12}=a,
\qquad
H_{21}=b.
\]
Then
\[
D_H
=
\begin{pmatrix}
-a&1\\
1&-b
\end{pmatrix},
\qquad
L_H
=
\begin{pmatrix}
a^2+1&-(a+b)\\
-(a+b)&b^2+1
\end{pmatrix}.
\]
Writing
\[
\tau=a^2+b^2+2,
\]
the eigenvalues are
\[
\lambda_\pm
=
\frac{1}{2}
\left(
\tau\pm\sqrt{\tau^2-4(ab-1)^2}
\right),
\]
and
\[
\det L_H=(ab-1)^2.
\]
Thus a nonzero exact shared core exists precisely when
\[
ab=1,
\]
which is exactly the trivial-holonomy condition. Away from exact consistency,
\[
\lambda_-\lambda_+=(ab-1)^2.
\]
When \(|ab-1|\) is small while \(\lambda_+\) stays bounded away from zero, the
soft workspace eigenvalue is quadratic in the cycle mismatch:
\[
\lambda_-
=
\frac{(ab-1)^2}{\lambda_+}.
\]
This example connects approximate spectral coherence directly to approximate
holonomy closure.

\topic{15}{Perturbation Stability Theorem}

The central robustness question is whether a selected low-energy subspace
persists when the estimated transfer maps change. The next lemma isolates the
finite-dimensional spectral fact needed for the answer.

\begin{lemma}[Separated spectral subspaces]\label{lem:separated-subspaces}
Let \(A\) and \(B=A+E\) be self-adjoint operators. Let \(P\) be an
\(A\)-spectral projector whose eigenvalues lie in \(( -\infty,a]\), and let
\(\widetilde Q\) be a \(B\)-spectral projector whose eigenvalues lie in
\([b,\infty)\), with \(b>a\). Then
\[
\|P\widetilde Q\|
\leq
\frac{\|E\|}{b-a}.
\]
\end{lemma}

\begin{proof}
On \(\ran P\) and \(\ran\widetilde Q\), set
\[
A_P=PAP|_{\ran P},
\qquad
B_Q=\widetilde QB\widetilde Q|_{\ran\widetilde Q}.
\]
The operator \(X=P\widetilde Q\) satisfies the Sylvester equation
\[
A_PX-XB_Q=-PE\widetilde Q.
\]
Let \(\mu=(a+b)/2\). Since
\[
\|e^{t(A_P-\mu I)}\|\leq e^{-(b-a)t/2},
\qquad
\|e^{-t(B_Q-\mu I)}\|\leq e^{-(b-a)t/2},
\]
the inverse Sylvester operator has the convergent integral representation
\[
X
=
\int_0^\infty
 e^{t(A_P-\mu I)}
 (PE\widetilde Q)
 e^{-t(B_Q-\mu I)}\,\dd t.
\]
Taking norms and integrating gives the result.
\end{proof}

\begin{theorem}[Spectral workspace stability]\label{thm:spectral-stability}
Let \(L\) and \(\widetilde L=L+E\) be self-adjoint positive semidefinite
operators on the same \(N\)-dimensional space. Suppose
\[
\lambda_d(L)<\lambda_{d+1}(L),
\qquad
g_d=\lambda_{d+1}(L)-\lambda_d(L)>0,
\]
and
\[
\eta=\|E\|<\frac{g_d}{2}.
\]
Let \(P_d\) and \(\widetilde P_d\) project onto the first \(d\) eigenvectors
of \(L\) and \(\widetilde L\), respectively. Then the rank-\(d\) cluster
remains separated and
\[
\boxed{
\|P_d-\widetilde P_d\|
\leq
\frac{\eta}{g_d-\eta}
\leq
\frac{2\eta}{g_d}.
}
\]
\end{theorem}

\begin{proof}
The min--max characterization gives the eigenvalue perturbation bound
\[
|\lambda_k(\widetilde L)-\lambda_k(L)|\leq\eta.
\]
Hence
\[
\lambda_{d+1}(\widetilde L)-\lambda_d(L)
\geq g_d-\eta
\]
and
\[
\lambda_{d+1}(L)-\lambda_d(\widetilde L)
\geq g_d-\eta.
\]
Apply Lemma~\ref{lem:separated-subspaces} first to
\(P_d(I-\widetilde P_d)\) and then, with the roles of \(L\) and
\(\widetilde L\) exchanged, to \((I-P_d)\widetilde P_d\). Orthogonal
projections of equal rank satisfy
\[
\|P_d-\widetilde P_d\|
=
\max\left\{
\|P_d(I-\widetilde P_d)\|,
\|(I-P_d)\widetilde P_d\|
\right\}.
\]
This yields the first bound. The second follows from \(\eta<g_d/2\).
\end{proof}

\subtopic{15.1 From transfer-map error to Laplacian error}

Suppose the edge maps are perturbed to
\[
\widetilde H_e=H_e+\Delta_e
\]
while the weights and vertex metrics remain fixed. Define
\[
\boxed{
\delta_H^2
=
\max_{i\in V}
\sum_{e:s(e)=i}w_e\|\Delta_e\|^2.
}
\]
Let \(F=D_{\widetilde H}-D_H\). Then
\[
(Fc)_e=-\sqrt{w_e}\,\Delta_ec_{s(e)},
\qquad
\|F\|\leq\delta_H.
\]
Since
\[
L_{\widetilde H}-L_H
=D_H^*F+F^*D_H+F^*F,
\]
we obtain
\[
\boxed{
\|L_{\widetilde H}-L_H\|
\leq
2\|D_H\|\delta_H+\delta_H^2.
}
\]

\begin{corollary}[Edgewise stability certificate]\label{cor:edge-stability}
If
\[
\eta_H=2\|D_H\|\delta_H+\delta_H^2
<\frac{g_d}{2},
\]
then
\[
\boxed{
\|P_d(H)-P_d(\widetilde H)\|
\leq
\frac{\eta_H}{g_d-\eta_H}.
}
\]
\end{corollary}

This converts declared edge-map uncertainty into an explicit bound on the
motion of the extracted workspace subspace.

\subtopic{15.2 Approximate operation closure}

Define the vertex and edge direct-sum representations
\[
\rho_V(a)=\bigoplus_{i\in V}\rho_i(a),
\qquad
\rho_E(a)=\bigoplus_{e\in E}\rho_{t(e)}(a).
\]
Lift the edgewise intertwining defects to
\[
(\mathfrak X_a c)_e
=
\sqrt{w_e}\,\Xi_e(a)c_{s(e)}.
\]
Then the exact commutator identity is
\[
\boxed{
D_H\rho_V(a)-\rho_E(a)D_H=-\mathfrak X_a.
}
\]
Furthermore,
\[
\|\mathfrak X_a\|
\leq
\beta(a),
\qquad
\beta(a)^2
=
\max_{i\in V}
\sum_{e:s(e)=i}w_e\|\Xi_e(a)\|^2.
\]

\begin{theorem}[Approximate operation-closure theorem]\label{thm:approx-operation-closure}
Let \(P_d\) project onto the first \(d\) eigenvectors of \(L_H\), set
\[
\lambda_d\leq\varepsilon,
\qquad
\mu=\lambda_{d+1}>0,
\]
and let \(Q_d=I-P_d\). Then every \(a\in\mathcal A\) satisfies
\[
\boxed{
\|Q_d\rho_V(a)P_d\|
\leq
\frac{
\|\rho_E(a)\|\sqrt{\varepsilon}+\beta(a)
}{\sqrt{\mu}}.
}
\]
\end{theorem}

\begin{proof}
For \(x\in\ran P_d\),
\[
\|D_Hx\|^2
=\langle x,L_Hx\rangle
\leq\varepsilon\|x\|^2.
\]
The commutator identity gives
\[
\|D_H\rho_V(a)x\|
\leq
\|\rho_E(a)\|\sqrt{\varepsilon}\|x\|
+
\beta(a)\|x\|.
\]
Because \(P_d\) and \(Q_d\) reduce \(L_H\), the two energy components are
orthogonal and
\[
\|D_HQ_dy\|\leq\|D_Hy\|.
\]
On \(\ran Q_d\), the spectrum of \(L_H\) is bounded below by \(\mu\), so
\[
\sqrt{\mu}\|Q_dy\|
\leq
\|D_HQ_dy\|.
\]
Apply these inequalities to \(y=\rho_V(a)x\), then take the supremum over unit
\(x\in\ran P_d\).
\end{proof}

\begin{corollary}[Recovery of exact operation closure]
If \(d=\dim\ker L_H\), \(\varepsilon=0\), and every edge intertwines \(a\)
exactly, then
\[
Q_d\rho_V(a)P_d=0.
\]
Thus Theorem~\ref{thm:approx-operation-closure} reduces to the exact
operation-closure result of Part~III.
\end{corollary}

\begin{corollary}[Closure under projector drift]\label{cor:closure-drift}
Let \(\widetilde P_d\) be another rank-\(d\) orthogonal projector and set
\[
\delta_P=\|P_d-\widetilde P_d\|.
\]
Then
\[
\boxed{
\|(I-\widetilde P_d)\rho_V(a)\widetilde P_d\|
\leq
\|Q_d\rho_V(a)P_d\|
+2\|\rho_V(a)\|\delta_P.
}
\]
Combining this estimate with Theorem~\ref{thm:spectral-stability} controls
operation leakage after noisy workspace extraction.
\end{corollary}

\begin{proof}
Subtract \(Q_d\rho_V(a)P_d\) and use
\[
I-\widetilde P_d-Q_d=P_d-\widetilde P_d.
\]
The two resulting perturbation terms are each bounded by
\(\|\rho_V(a)\|\delta_P\).
\end{proof}

\subtopic{15.3 Core--residual path accumulation}

Let a path have vertices
\[
i_0\xrightarrow{e_1}i_1\xrightarrow{e_2}\cdots
\xrightarrow{e_m}i_m.
\]
Write
\[
T_r=T_{e_r},
\quad
H_r=H_{e_r},
\quad
V_r=V_{e_r},
\]
and let
\[
\iota_i^C:C_i\hookrightarrow X_i,
\qquad
\iota_i^R:R_i\hookrightarrow X_i
\]
be the canonical inclusions. Empty products below are identities.

\begin{theorem}[Pathwise core--residual Duhamel formula]\label{thm:residual-duhamel}
The deviation of full path transport from pure core transport is exactly
\[
\boxed{
\begin{aligned}
&T_m\cdots T_1\iota_{i_0}^C
-
\iota_{i_m}^C H_m\cdots H_1\\
&\qquad=
\sum_{r=1}^{m}
T_m\cdots T_{r+1}
\iota_{i_r}^R V_r
H_{r-1}\cdots H_1.
\end{aligned}
}
\]
Consequently,
\[
\boxed{
\begin{aligned}
\|Q_{i_m}T_m\cdots T_1P_{i_0}\|
\leq
\sum_{r=1}^{m}
&\left(\prod_{\ell=r+1}^{m}\|T_\ell\|\right)
\|V_r\|\\
&\times
\left(\prod_{\ell=1}^{r-1}\|H_\ell\|\right),
\end{aligned}
}
\]
and the same right-hand side bounds
\[
\|P_{i_m}T_m\cdots T_1P_{i_0}-H_m\cdots H_1\|.
\]
\end{theorem}

\begin{proof}
For every step,
\[
T_r\iota_{i_{r-1}}^C
=
\iota_{i_r}^CH_r+
\iota_{i_r}^RV_r.
\]
Telescope the sequence
\[
T_m\cdots T_{r+1}\iota_{i_r}^C H_r\cdots H_1,
\qquad
r=0,1,\ldots,m.
\]
The difference between consecutive terms is exactly the \(r\)-th summand in
the displayed formula. Projection onto \(R_{i_m}\) or \(C_{i_m}\), followed
by submultiplicativity of the operator norm, gives the two bounds.
\end{proof}

\begin{corollary}[Uniform path bound]\label{cor:uniform-path-bound}
If along the path
\[
\|T_r\|\leq\tau,
\qquad
\|H_r\|\leq h,
\qquad
\|V_r\|\leq\nu,
\]
then the leakage and core-transport defects are bounded by
\[
\boxed{
\nu\sum_{r=1}^{m}\tau^{m-r}h^{r-1}
=
\begin{cases}
\displaystyle
\nu\frac{\tau^m-h^m}{\tau-h},&\tau\neq h,\\[1.1em]
\displaystyle
m\nu\tau^{m-1},&\tau=h.
\end{cases}
}
\]
In particular, if \(\tau\leq1\) and \(h\leq1\), the bound is at most
\(m\nu\).
\end{corollary}

The linear factor \(m\) in the merely nonexpansive regime is important. Small
single-edge leakage does not by itself justify arbitrarily long transfer
chains. Long-horizon stability requires contraction, cancellation, error
correction, or leakage that decreases with path length.

\subtopic{15.4 Residual-mediated reentry}

For two consecutive transfers,
\[
P_{i_2}T_2T_1P_{i_0}
=
H_2H_1+U_2V_1,
\]
\[
Q_{i_2}T_2T_1P_{i_0}
=
V_2H_1+K_2V_1.
\]
Thus leakage out of the core is first order in the blocks \(V_r\), while an
excursion that leaves and then returns to the core contains at least one
\(V\)-block and one later \(U\)-block. For a fixed path length, if
\[
\|U_r\|,\|V_r\|=O(\delta)
\]
and the diagonal blocks remain bounded, the residual-mediated correction to
pure core-to-core transport is \(O(\delta^2)\), whereas terminal residual
leakage can be \(O(\delta)\). This distinction separates loss of core
invariance from distortion caused by temporary excursions through local
residual degrees of freedom.
